% PosTextuais/ApendiceA.tex
\begin{apendicesenv}
\partapendices

\chapter{FORMULAÇÃO MATEMÁTICA COMPLETA E LINEARIZADA DO SOLVER MILP}
\label{apendice:a}

O problema de alocação ótima e escalonamento de manutenções com restrição de capacidade e orçamento para a frota do Estado de Santa Catarina é formulado canonicamente como segue:

\begin{equation}
    \min \sum_{t=1}^T \sum_{i=1}^N \left[ \alpha \cdot c_{it} x_{it} + \beta \cdot r_{it} (1 - x_{it}) \right]
\end{equation}

\noindent Restrições Operacionais e Orçamentárias:
\begin{enumerate}
    \item Teto Orçamentário Mensal: $\sum_{i=1}^N c_{it} x_{it} \le B_t, \quad \forall t \in \{1, \dots, T\}$
    \item Capacidade Operacional de Oficinas: $\sum_{i=1}^N h_{it} x_{it} \le H_t, \quad \forall t \in \{1, \dots, T\}$
    \item Manutenções Mandatórias de Segurança: $x_{it} \ge m_{it}, \quad \forall (i, t) \mid m_{it} = 1$
    \item Integralidade e Domínio: $x_{it} \in \{0, 1\}, \quad \forall i, t$
\end{enumerate}

\noindent\textbf{Métricas Poliedrais e Complexidade Computacional da Instância S0:}

A instância canônica S0 contempla 72.000 variáveis binárias de decisão, correspondentes a 2.000 veículos avaliados ao longo de 36 períodos mensais. O modelo apresenta 6.072 restrições lineares reportadas, compreendendo limites orçamentários, capacidades operacionais e exigências de manutenção mandatória. A quantidade de restrições efetivamente ativas na solução ótima deverá ser determinada a partir dos resultados do solver, não podendo ser inferida exclusivamente da formulação matemática.
